\section{Анализ предметной области}

\subsection{Понятие рекомендательных алгоритмов}
Пусть рассматриваются предметы, представленные какими-либо данными и пользователи, которые их используют.
Пользователи присваивают предметам оценки в соответствии со своими субъективными предпочтениями.
Оценкой является число из определённого множества или отрезка.
Считается, что чем выше оценка, тем предмет субъективно нужнее пользователю.
Далее рассмотрена одна из постановок цели рекомендательной системы.
Необходимо предсказать оценки, которые выбранный пользователь присвоит новым для него предметам, после чего сформировать список предметов, которые данный пользователь определит как субъективно нужные~\cite{aggarwal}.
Таким образом, объект исследования -- рекомендательные системы -- это программы, которые информируют пользователя о субъективно нужных для него объектах.

\subsection[Формализация задачи рекомендательных\\систем]{Формализация задачи рекомендательных систем}
Рассматривается постановка задачи рекомендательных алгоритмов.
Пусть:
\begin{itemize}
    \item $U$ -- множество пользователей программы;
    \item $X$ -- множество из $M$ объектов, с которыми взаимодействуют пользователи;
    \item $W_i \subset X$ -- множество объектов, которые оценил пользователь $u_i$, где $i \in \overline{1, N}$, $N$ -- количество пользователей;
    \item $\alpha_{i,j}$ -- оценка объекта $X_j$ пользователем $U_i$, где $X_j \in W_i, i \in \overline{1, N}$.
\end{itemize}
Матрицей оценок называется матрица $Y$ размером $N \times M$, где элемент $y_{i,j}$ строки $i$ и столбца $j$ содержит либо оценку пользователя $U_i$ для объекта $X_j$, либо метку о её отсутствии~\cite{aggarwal}.
Также матрица оценок может содержать неявные оценки объектов пользователем, принимающие значения $0$ и $1$.
Такие оценки собираются автоматически во время работы рекомендательной системы.
Входными данными является следующее:
\begin{itemize}
    \item матрица оценок;
    \item матрица неявных оценок;
    \item данные об объектах.
\end{itemize}
Необходимо выбрать из множества $X \setminus W_i$ один или несколько объектов таким образом, чтобы его оценка пользователем $U_i$ была максимальна~\cite{aggarwal}.
Таким образом, выходными данными является список объектов $x_i \in X$, отсортированный по убыванию предполагаемой оценки от пользователя.